% chapter2.tex — บทที่ 2 เอกสารและงานวิจัยที่เกี่ยวข้อง
\chapter{เอกสารและงานวิจัยที่เกี่ยวข้อง}

\section{แนวคิด ทฤษฎีที่เกี่ยวข้อง}
(สังเคราะห์แนวคิด/ทฤษฎีหลักที่เป็นกรอบของงานวิจัย พร้อมอ้างอิงแหล่งที่มา เช่น \parencite{example2024})

\section{งานวิจัยที่เกี่ยวข้อง}
(สรุปงานวิจัยในและต่างประเทศที่เกี่ยวข้อง เรียงตามลำดับเวลาหรือประเด็น)

\begin{longtable}{|p{3.2cm}|p{3.4cm}|p{3.8cm}|p{4cm}|}
\caption{สรุปงานวิจัยที่เกี่ยวข้อง (ตัวอย่างรูปแบบตาราง)} \label{tab:2-1}\\
\hline
\tblhead{ผู้วิจัย/ปี} & \tblhead{วัตถุประสงค์} & \tblhead{วิธีดำเนินการ} & \tblhead{ผลการวิจัยที่สำคัญ} \\
\hline
\endfirsthead
\multicolumn{4}{c}{\tablename\ \thetable\ (ต่อ)}\\
\hline
\tblhead{ผู้วิจัย/ปี} & \tblhead{วัตถุประสงค์} & \tblhead{วิธีดำเนินการ} & \tblhead{ผลการวิจัยที่สำคัญ} \\
\hline
\endhead
\hline
\endfoot
\hline
\endlastfoot
[ผู้แต่ง, ปีที่พิมพ์] & [วัตถุประสงค์โดยย่อ] & [กลุ่มตัวอย่าง/เครื่องมือ] & [ผลสรุปที่เกี่ยวข้อง] \\
\hline
[ผู้แต่ง, ปีที่พิมพ์] & [วัตถุประสงค์โดยย่อ] & [กลุ่มตัวอย่าง/เครื่องมือ] & [ผลสรุปที่เกี่ยวข้อง] \\
\hline
\end{longtable}

\section{กรอบแนวคิดในการวิจัย}
(นำเสนอกรอบแนวคิด (Conceptual Framework) เชื่อมโยงตัวแปรต้น-ตาม อาจแทรกภาพประกอบ เช่น)

% ตัวอย่างการแทรกภาพประกอบ (ต้องอัปโหลดไฟล์ภาพไปที่โฟลเดอร์ images/)
% \begin{figure}[h]
%   \centering
%   \includegraphics[width=0.7\textwidth]{images/framework.png}
%   \caption{กรอบแนวคิดในการวิจัย}
%   \label{fig:2-1}
% \end{figure}
